\chapter{Réalisation}

\section*{Introduction}
Ce chapitre décrit la mise en œuvre de la conception : les paquets ROS~2 ajoutés ou
modifiés, la passerelle de navigation, les composants du dashboard, l'outillage de
lancement et les tests automatisés.

\section{Organisation des fichiers}
Le travail s'insère dans le dépôt existant. Le tableau~\ref{tab:fichiers} liste les
principaux fichiers ajoutés.

\begin{table}[!htb]
\centering\footnotesize
\begin{tabularx}{\linewidth}{@{}>{\ttfamily\raggedright\arraybackslash}p{7.4cm}X@{}}
\toprule
\normalfont\textbf{Fichier} & \textbf{Rôle} \\
\midrule
navigation\_msgs/srv/ SubmitNavigationGoal.srv & Interface du service de soumission. \\
navigation/scripts/ dashboard\_navigation\_node.py & Passerelle : services, action Nav2, TF, état JSON. \\
navigation/scripts/ dashboard\_navigation\_math.py & Conversions de repères et lecture de la costmap, sans dépendance ROS. \\
navigation/launch/ dashboard\_navigation.launch.py & Lancement de la passerelle et, en option, de rosbridge. \\
navigation/test/ test\_dashboard\_navigation\_math.py & Tests unitaires Python des conversions. \\
frontend/src/hooks/ useNavigation.js & Abonnement à l'état, fraîcheur, appels de services. \\
frontend/src/utils/ navigation.js & Validation du protocole et conditions d'envoi. \\
frontend/src/components/ NavigationPanel.jsx & Panneau de destination, d'état et d'annulation. \\
frontend/tests/ navigation.test.mjs & Tests unitaires JavaScript. \\
scripts/run-all.sh & Lancement de toute la chaîne en une commande. \\
\bottomrule
\end{tabularx}
\caption{Principaux fichiers ajoutés.}
\label{tab:fichiers}
\end{table}

Les fichiers existants modifiés sont principalement les \code{CMakeLists.txt} et
\code{package.xml} des paquets \code{navigation} et \code{navigation\_msgs} (génération de
l'interface, installation des scripts, dépendances), la page \code{Map.jsx}, les
composants de carte 2D et 3D, le hook de connexion ROS et la page des paramètres.

\section{Passerelle ROS~2}

\subsection{Paramètres}
La passerelle est un nœud Python (\code{rclpy}). Ses paramètres sont donnés dans le
tableau~\ref{tab:params}.

\begin{table}[htbp]
\centering\small
\begin{tabular}{@{}lll@{}}
\toprule
\textbf{Paramètre} & \textbf{Défaut} & \textbf{Rôle} \\
\midrule
\code{world\_id} & -- & Monde du dashboard correspondant à la simulation \\
\code{base\_frame} & \code{base\_link} & Repère du robot lu dans TF \\
\code{world\_offset\_x}, \code{\_y} & 0,0 & Position de l'origine de \code{map} dans le monde (m) \\
\code{world\_offset\_yaw} & 0,0 & Rotation de \code{map} dans le monde (rad) \\
\code{occupied\_threshold} & 65 & Seuil de coût d'une cellule occupée \\
\code{use\_sim\_time} & \code{true} & Imposé par le fichier de lancement \\
\bottomrule
\end{tabular}
\caption{Paramètres de la passerelle.}
\label{tab:params}
\end{table}

\subsection{Fonctionnement}
Le nœud s'abonne à la costmap globale avec une qualité de service \emph{transient local},
identique à celle de l'éditeur Nav2, afin de recevoir la dernière grille même s'il démarre
après Nav2. Un temporisateur de 0,25~s calcule à chaque cycle la pose du robot et la
disponibilité, puis publie l'état. Les principaux traitements sont :
\begin{itemize}
  \item \textbf{soumission} : contrôle de l'identifiant, du but actif, du monde, des
    coordonnées, de la disponibilité, conversion par l'équation~\eqref{eq:world2map}, contrôle
    de la cellule, puis envoi asynchrone du but à Nav2 ;
  \item \textbf{acceptation} : à la réponse de Nav2, passage en \code{navigating} ou en
    \code{rejected}, et envoi d'une annulation mémorisée le cas échéant ;
  \item \textbf{retour} : mise à jour de la distance restante ;
  \item \textbf{résultat} : conversion du statut Nav2 en \code{succeeded},
    \code{canceled} ou \code{failed} ;
  \item \textbf{annulation} : passage en \code{canceling} et demande d'annulation à Nav2 ;
    la phase terminale n'est affichée qu'à réception du résultat.
\end{itemize}
Les calculs de repères et de cellules sont isolés dans un module Python sans dépendance
ROS, ce qui permet de les tester sans lancer de simulation.

\subsection{Fichier de lancement}
Le fichier \code{dashboard\_navigation.launch.py} démarre la passerelle avec
\code{use\_sim\_time} et inclut le serveur rosbridge, sauf si l'argument
\code{start\_rosbridge:=false} est donné. Les décalages de repère sont passés en arguments :
\begin{lstlisting}
ros2 launch navigation dashboard_navigation.launch.py \
  world_id:=warehouse-harmonic \
  world_offset_x:=-1.2639 world_offset_y:=-3.9049 world_offset_yaw:=0.0
\end{lstlisting}

\section{Interface web}

\subsection{Logique de navigation}
Le hook React \code{useNavigation} s'abonne à l'état publié par la passerelle au moyen de
\code{roslib}. Chaque message est validé par \code{parseNavigationState} : version,
types des champs, phase connue et pose finie. Un message non conforme n'est pas affiché
comme une donnée réelle. L'âge du dernier message est réévalué toutes les 0,5~s ; au-delà
de 2,5~s, l'état est considéré périmé.

La fonction \code{navigationBlockReason} centralise les conditions d'envoi et retourne le
motif de blocage à afficher (mode démonstration, déconnexion, données périmées, mauvais
monde, passerelle non prête, but actif). Elle est réévaluée au moment exact du clic, et
non seulement au dernier rendu. Un seul appel de service peut être en attente ; en
l'absence de réponse sous 6~s, l'interface signale que le résultat est incertain et ne
renvoie pas la commande.

\subsection{Panneau de navigation et carte}
Le composant \code{NavigationPanel} affiche la phase et son
détail, la distance restante, le motif de blocage éventuel et les erreurs. L'opérateur
choisit une destination en cliquant sur la carte 2D après avoir activé \emph{Choose on
map}, ou en saisissant X et Y en mètres et le cap en degrés. Les boutons \emph{Go to
destination} et \emph{Cancel navigation} ne sont actifs que lorsque l'action est permise.
Un message rappelle que fermer la page n'annule pas la navigation.

La page \emph{Live map} a été adaptée (figure~\ref{fig:livemap}) :
\begin{itemize}
  \item le robot est dessiné à partir de la pose ROS uniquement lorsque l'état est frais
    et que le monde correspond ; la pose de la base de données n'est plus utilisée en
    remplacement ;
  \item un marqueur indique la destination choisie ;
  \item le mode démonstration est signalé explicitement et une erreur de l'API n'active
    plus implicitement les robots de démonstration ;
  \item la projection du robot et du trajet sur le plan 2D utilise la même transformation
    que le monde SDF, et l'orientation ROS est convertie pour le repère SVG ;
  \item la vue 3D n'applique plus de décalage arbitraire et oriente le robot selon sa pose ;
    le robot et le trajet décoratifs ont été supprimés.
\end{itemize}
\begin{figure}[htbp]
\centering
\capture{live_map_navigation.png}{Page Live map.}
\caption{Page \emph{Live map} : robot \code{amr\_x} affiché à sa position ROS sur le plan
du monde \code{warehouse\_harmonic}.}
\label{fig:livemap}
\end{figure}

Enfin, la connexion rosbridge se rétablit automatiquement après une fermeture du
WebSocket, et une nouvelle adresse saisie dans \emph{Settings} est appliquée sans
recharger la page.

\section{Lancement de la chaîne complète}
La chaîne complète comporte cinq processus : simulation, Nav2, passerelle avec rosbridge,
backend et frontend. Pour la rendre reproductible, un script \code{scripts/run-all.sh} les
démarre dans l'ordre, attend que chacun soit prêt avant de lancer le suivant, et arrête
l'ensemble par un simple \emph{Ctrl+C} (figure~\ref{fig:lancement}). Il envoie également la
pose initiale à AMCL, ce qui évite une intervention manuelle dans RViz.

\begin{lstlisting}
cd amr
bash scripts/run-all.sh          # fenetre Gazebo ouverte
bash scripts/run-all.sh rviz     # avec RViz en plus
\end{lstlisting}

\begin{figure}[htbp]
\centering
\resizebox{\linewidth}{!}{%
\begin{tikzpicture}[node distance=0.55cm]
  \node[sim] (a) {1. Gazebo\\\scriptsize attente \code{/clock}};
  \node[ros,right=of a] (b) {2. Nav2\\\scriptsize attente AMCL actif};
  \node[ros,right=of b] (c) {pose initiale\\\scriptsize \code{/initialpose}};
  \node[neuf,right=of c] (d) {3. Passerelle\\\scriptsize attente \code{ready}};
  \node[web,right=of d] (e) {4. Backend\\\scriptsize :8000};
  \node[web,right=of e] (f) {5. Frontend\\\scriptsize :5173};
  \draw[fl] (a) -- (b); \draw[fl] (b) -- (c); \draw[fl] (c) -- (d);
  \draw[fl] (d) -- (e); \draw[fl] (e) -- (f);
\end{tikzpicture}}
\caption{Séquence de démarrage du script \code{run-all.sh}.}
\label{fig:lancement}
\end{figure}

\section{Tests automatisés}
Deux jeux de tests unitaires ne nécessitent ni ROS ni Gazebo :
\begin{itemize}
  \item \textbf{Python} (5 tests) : aller-retour entre les repères, repère carte tourné,
    cellules occupées, inconnues et hors carte, grille d'origine tournée, grille invalide ;
  \item \textbf{JavaScript} (5 tests) : conditions d'envoi, blocage des buts concurrents,
    rejet des états mal formés, validité des coordonnées nulles ou négatives,
    unicité et format des identifiants de requête.
\end{itemize}
Un test navigateur optionnel, fondé sur Playwright et un faux serveur rosbridge, vérifie
en plus le parcours complet de l'interface (envoi, conversion degrés/radians, annulation,
sélection sur la carte, blocage d'un mauvais monde). Il ne remplace pas les essais dans
Gazebo présentés au chapitre suivant.

\section*{Conclusion}
La passerelle, l'interface web et l'outillage de lancement ont été réalisés et intégrés
à l'environnement du projet. Les tests unitaires couvrent les conversions de repères et la logique de
l'interface. Il reste à vérifier le comportement de l'ensemble dans la simulation, ce qui
fait l'objet du chapitre suivant.
